\chapter{The system at a glance}

Before diving into any single file, let's see the whole machine. This codebase takes a
$672\times672$ RGB image and, in one forward pass, predicts three things for every
object it finds:

\begin{enumerate}[leftmargin=1.4em,itemsep=1pt]
  \item a \textbf{bounding box} and a \textbf{class} (this is ordinary object detection);
  \item a \textbf{polygon outline} of the object (segmentation);
  \item the \textbf{distance} to the object in meters.
\end{enumerate}

It is a faithful TensorFlow~2.16 re-implementation of Ultralytics YOLOv8, extended with
two extra prediction ``heads''. The assembly happens in \file{models/yolo\_v8.py}
(\code{build\_yolov8}).

\section{The end-to-end path}
Every image flows through four stages:

\begin{center}
\begin{tikzpicture}[node distance=7mm and 10mm]
  \node[blkn] (img) {Input image\\$672\times672\times3$};
  \node[blk,right=of img] (bb) {\textbf{Backbone}\\CSPDarkNet-V8-S\\\scriptsize feature extractor};
  \node[blk,right=of bb] (dec) {\textbf{Decoder}\\FPN-PAN\\\scriptsize multi-scale fusion};
  \node[blkg,right=of dec] (head) {\textbf{Heads}\\6 branches\\\scriptsize per pixel};
  \node[blka,below=9mm of head] (nms) {\textbf{Detection generator}\\NMS + decode\\\scriptsize (inference only)};
  \node[blkn,left=of nms] (out) {Final\\detections};
  \draw[ar] (img)--(bb); \draw[ar] (bb)--(dec); \draw[ar] (dec)--(head);
  \draw[ar] (head)--(nms); \draw[ar] (nms)--(out);
  \node[lbl,below=0.5mm of bb] {strides 8/16/32};
  \node[lbl,below=0.5mm of dec] {3 feature maps};
\end{tikzpicture}
\end{center}

\begin{teacher}
A \textbf{backbone} is a stack of convolutional layers that turns raw pixels into
compact ``feature maps'' --- grids of numbers where each cell summarizes a patch of the
image. A \textbf{decoder} (here a feature pyramid) mixes information across resolutions
so the model can see both small and large objects. The \textbf{heads} are the final
thin layers that read those features and output the actual predictions. We unpack each
of these in the architecture chapter.
\end{teacher}

\section{Three tiers, one codebase}
The same code supports three configurations, chosen by an experiment YAML file. They
differ only in which heads are turned on:

\begin{center}
\renewcommand{\arraystretch}{1.3}
\begin{tabular}{>{\raggedright\arraybackslash}p{4.5cm} l l}
\toprule
\textbf{Config (\file{configs/experiments/yolo/})} & \textbf{Heads} & \textbf{Task} \\
\midrule
\file{yolov8\_bbox.yaml}      & box + cls & detection only \\
\file{yolov8\_poly.yaml}      & box + cls + polygon & detection + segmentation \\
\file{yolov8\_poly\_dist.yaml} & all 6 heads & detection + segmentation + distance \\
\bottomrule
\end{tabular}
\end{center}

\noindent The full tier, \file{yolov8\_poly\_dist.yaml}, is the authoritative reference
configuration --- when in doubt about a hyperparameter, that file is the source of truth.

\section{The six heads}
Every head produces a prediction at \emph{every pixel} of \emph{every} one of the three
feature-map scales. Here is what each one outputs.

\begin{center}
\renewcommand{\arraystretch}{1.3}
\begin{tabular}{l c >{\raggedright\arraybackslash}p{8.2cm}}
\toprule
\textbf{Head} & \textbf{Channels} & \textbf{Meaning} \\
\midrule
\code{box}        & 64 & box position, as a ``distribution'' --- 4 sides $\times$ 16 bins (DFL) \\
\code{cls}        & 39 & one logit per class \\
\code{poly\_angle} & 24 & per-vertex sub-bin angle offset (one per polygon bin) \\
\code{poly\_dist}  & 24 & per-vertex radial distance (one per polygon bin) \\
\code{poly\_conf}  & 24 & per-vertex ``is there a vertex here?'' gate \\
\code{dist}       & 1  & object distance, on a log scale \\
\bottomrule
\end{tabular}
\end{center}

\begin{hood}
The heads live in \file{models/head.py} (\code{YoloV8Head}). The three polygon heads
(\code{poly\_angle}, \code{poly\_dist}, \code{poly\_conf}) together encode a 24-vertex
outline in the ``radial'' PolyYOLO format --- the subject of its own chapter. The
\code{box} head's 64 channels are $4\times16$: each of the four box sides is predicted
as a little probability distribution over 16 distance bins, a trick called
\emph{Distribution Focal Loss}. The heads are pinned to \code{float32} even when the
rest of the model runs in \code{mixed\_bfloat16}.
\end{hood}

\section{What ``training'' produces}
During training the model runs in a mode where the heads emit \emph{raw} numbers and a
loss compares them to ground-truth labels. During inference (\code{deploy=True}) a final
stage --- the \textbf{detection generator} (\file{models/detection\_generator.py}) ---
decodes those raw numbers into clean boxes, runs \textbf{non-maximum suppression} to
remove duplicates, and converts the polygon and distance outputs into human-readable
form.

\begin{teacher}
\textbf{Non-maximum suppression (NMS)} is a clean-up step. A detector typically fires on
the same object from several neighboring pixels, producing a cluster of overlapping
boxes. NMS keeps the highest-scoring box in each cluster and suppresses the rest, so each
object is reported once.
\end{teacher}

With the map in hand, the next chapter builds up the detection ideas themselves --- so
that when we open the code, every term already means something.
